% ============================================================
\chapter*{Preface}
\addcontentsline{toc}{chapter}{Preface}
% ============================================================

Riemannian geometry stands as one of the fundamental pillars of modern mathematics,
lying at the crossroads of analysis, algebra, and topology. Born from the visionary
work of Bernhard Riemann in his celebrated 1854 \emph{Habilitationsschrift}, entitled
\emph{\"Uber die Hypothesen, welche der Geometrie zu Grunde liegen}, this discipline
has undergone spectacular development, culminating in results as profound as
Perelman's proof of the Poincar\'e conjecture via Ricci flow.

\section*{Goals of this course}

This course is aimed at Master's and PhD students in pure mathematics. It provides
a rigorous and thorough presentation of the foundations of Riemannian geometry,
from the definition of Riemannian manifolds through comparison theorems and an
introduction to Ricci flow.

The main objectives are:
\begin{itemize}
    \item Master the formalism of connections, covariant differentiation, and
          parallel transport.
    \item Understand in depth the Levi-Civita connection and its characterizing
          properties.
    \item Study geodesics, the exponential map, and minimality properties.
    \item Develop the tensor calculus of curvature: the Riemann tensor, sectional
          curvature, Ricci curvature, and scalar curvature.
    \item Establish the major global theorems: Bonnet--Myers, Hadamard--Cartan,
          and the comparison theorems of Rauch and Toponogov.
    \item Explore the geometry of symmetric spaces and submanifolds.
    \item Provide an introduction to Ricci flow, motivated by the Poincar\'e
          conjecture.
\end{itemize}

\section*{Prerequisites}

The reader is assumed to be familiar with:
\begin{itemize}
    \item The theory of smooth manifolds: atlases, charts, diffeomorphisms,
          tangent bundles, differential forms.
    \item Multilinear algebra and basic tensor calculus: tensor products,
          exterior algebra, contraction.
    \item Elements of general topology: compactness, connectedness, covering spaces.
    \item Ordinary differential equations: the Cauchy--Lipschitz (Picard--Lindel\"of)
          theorem, smooth dependence on initial conditions.
    \item Basics of Lie group theory and Lie algebras.
\end{itemize}

Prior familiarity with the differential geometry of curves and surfaces in $\R^3$
is helpful but not essential; the relevant concepts will be revisited in the
intrinsic framework.

\section*{Organization of the course}

The course is organized into eleven chapters, structured as follows:

\begin{description}
    \item[Chapter 1 --- Riemannian Manifolds.]
    We introduce Riemannian metrics, the associated metric space structure, and
    the fundamental examples: spheres $S^n$, hyperbolic spaces $\mathbb{H}^n$,
    flat tori, and Lie groups with bi-invariant metrics. Isometries and completeness
    are discussed.

    \item[Chapter 2 --- Connections and Covariant Derivatives.]
    We develop the general theory of connections on vector bundles: axiomatic
    definition, parallel transport, torsion, and holonomy.

    \item[Chapter 3 --- The Levi-Civita Connection.]
    We prove the fundamental theorem of Riemannian geometry: existence and
    uniqueness of the torsion-free, metric-compatible connection. Christoffel
    symbols are computed explicitly on examples.

    \item[Chapter 4 --- Geodesics and the Exponential Map.]
    We study geodesics as locally length-minimizing curves, the exponential map,
    the Gauss lemma, normal coordinates, and the injectivity radius.

    \item[Chapter 5 --- Curvature: The Riemann Tensor.]
    We introduce the Riemann curvature tensor via the non-commutativity of
    covariant differentiation, its fundamental symmetries, and the Bianchi
    identities.

    \item[Chapter 6 --- Sectional, Ricci, and Scalar Curvature.]
    We define the various contractions of the Riemann tensor and compute them
    on the model spaces.

    \item[Chapter 7 --- Comparison Theorems.]
    We present the theorems of Rauch, Toponogov, and Bishop--Gromov, relating
    curvature bounds to global geometry.

    \item[Chapter 8 --- Symmetric Spaces.]
    We study Riemannian symmetric spaces in the sense of Cartan, their connection
    to Lie groups, and the classification into compact and noncompact types.

    \item[Chapter 9 --- Bonnet--Myers and Hadamard--Cartan.]
    We prove these two fundamental theorems illustrating the influence of
    curvature sign on global topology.

    \item[Chapter 10 --- Submanifold Geometry.]
    We develop the theory of submanifolds: the second fundamental form,
    the Gauss--Codazzi--Ricci equations, hypersurfaces, and mean curvature.

    \item[Chapter 11 --- Introduction to Ricci Flow.]
    We present Hamilton's Ricci flow, local existence, self-similar solutions
    (solitons), and an overview of Perelman's strategy.
\end{description}

\section*{Conventions and notation}

Throughout this course, we adopt the following conventions:
\begin{itemize}
    \item Manifolds are smooth ($C^\infty$), Hausdorff, second-countable, and
          connected unless otherwise stated.
    \item $(M, g)$ denotes a Riemannian manifold, $\nabla$ the Levi-Civita
          connection, $R$ the curvature tensor.
    \item Sign convention for the Riemann tensor:
          \[
              R(X,Y)Z = \nabla_X \nabla_Y Z - \nabla_Y \nabla_X Z - \nabla_{[X,Y]} Z.
          \]
    \item The Einstein summation convention is used: a repeated index in
          contravariant and covariant positions is implicitly summed.
    \item $\mathfrak{X}(M)$ denotes the space of smooth vector fields on $M$.
    \item $\Gamma(E)$ denotes the space of smooth sections of a vector bundle $E$.
    \item $T_pM$ denotes the tangent space at $p$, $T^*_pM$ the cotangent space.
    \item $\inner{\cdot}{\cdot}$ or $g(\cdot,\cdot)$ denotes the Riemannian inner product.
    \item $\norm{v} = \sqrt{\inner{v}{v}}$ denotes the Riemannian norm.
    \item $d(p,q)$ denotes the Riemannian distance between $p$ and $q$.
    \item $B(p,r) = \{q \in M : d(p,q) < r\}$ denotes the geodesic ball.
    \item $\operatorname{Ric}$ denotes the Ricci tensor, $S$ or $\operatorname{Scal}$
          the scalar curvature.
    \item $K(\sigma)$ denotes the sectional curvature of the $2$-plane $\sigma$.
\end{itemize}

\begin{keyformulas}[title=\textbf{Reference model spaces}]
The three families of model spaces of constant curvature:
\begin{enumerate}
    \item \textbf{Sphere} $S^n(r)$: sectional curvature $K = 1/r^2$,
          $\pi_1 = 0$ for $n \geq 2$.
    \item \textbf{Euclidean space} $\R^n$: curvature $K = 0$, flat.
    \item \textbf{Hyperbolic space} $\mathbb{H}^n(r)$: curvature $K = -1/r^2$,
          simply connected, noncompact.
\end{enumerate}
\end{keyformulas}

\section*{Guiding theme: curvature and topology}

The central theme of this course is the relationship between \emph{curvature}
(a local, analytic datum) and \emph{topology} (a global, qualitative datum).
This interplay is manifested in the following results:

\begin{center}
\renewcommand{\arraystretch}{1.3}
\begin{tabular}{lll}
\hline
\textbf{Curvature condition} & \textbf{Topological consequence} & \textbf{Theorem} \\
\hline
$\operatorname{Ric} \geq (n-1)\kappa > 0$ & $M$ compact, $\pi_1$ finite & Bonnet--Myers \\
$K \leq 0$ everywhere & $\widetilde{M} \cong \R^n$ & Hadamard--Cartan \\
$K \geq \delta > 0$ & $\operatorname{diam}(M) \leq \pi/\sqrt{\delta}$ & Bonnet--Myers \\
$K > 1/4$-pinched & $M$ homeomorphic to $S^n$ & Sphere theorem \\
\hline
\end{tabular}
\end{center}

\section*{Study recommendations}

Riemannian geometry is a discipline where geometric intuition and analytical rigor
must complement each other. We recommend:
\begin{itemize}
    \item Always checking general formulas on concrete examples ($S^2$, $\mathbb{H}^2$,
          flat torus).
    \item Drawing geometric situations, even schematically.
    \item Solving the proposed exercises, which are an integral part of the course.
    \item Systematically comparing intrinsic and coordinate-based approaches.
\end{itemize}

\section*{References}

\begin{enumerate}
    \item M.~P.~do Carmo, \emph{Riemannian Geometry}, Birkh\"auser, 1992.
    \item J.~M.~Lee, \emph{Riemannian Manifolds: An Introduction to Curvature},
          Springer, 1997.
    \item S.~Gallot, D.~Hulin, J.~Lafontaine, \emph{Riemannian Geometry},
          Universitext, Springer, 2004.
    \item P.~Petersen, \emph{Riemannian Geometry}, GTM 171, Springer, 2016.
    \item J.~Jost, \emph{Riemannian Geometry and Geometric Analysis}, Springer, 2017.
    \item J.~Cheeger, D.~Ebin, \emph{Comparison Theorems in Riemannian Geometry},
          AMS Chelsea, 2008.
    \item S.~Kobayashi, K.~Nomizu, \emph{Foundations of Differential Geometry},
          vols.~I and II, Wiley, 1963--1969.
    \item B.~Chow, P.~Lu, L.~Ni, \emph{Hamilton's Ricci Flow}, AMS, 2006.
    \item M.~Berger, \emph{A Panoramic View of Riemannian Geometry}, Springer, 2003.
\end{enumerate}

\bigskip
\hfill\emph{The author, March 2026}
